\documentclass[10pt]{mypackage}

\usepackage{mlmodern}
%\usepackage{newpxtext,eulerpx,eucal}
%\renewcommand*{\mathbb}[1]{\varmathbb{#1}}

\usepackage{homework}
%\usepackage{notes}

%\usepackage[ backend=bibtex, style = alphabetic, sorting=ynt ]{biblatex}
%\addbibresource{  }

\usepackage{parskip}

\fancyhf{}
\fancyhead[R]{Avinash Iyer}
\fancyhead[L]{Algebraic Topology: Homework 18}
\fancyfoot[C]{\thepage}

\setcounter{secnumdepth}{0}

\begin{document}
\RaggedRight
\begin{problem}[Problem 1]
  Prove that the only covering space $ \mathbb{RP}^2\rightarrow X $ is the identity map on $ \mathbb{RP}^2 $.
\end{problem}
\begin{solution}
  Under the assumption that $X$ is a (necessarily two-dimensional) CW complex, we observe that if $e^{0}_{\alpha}\in X$ is any $0$-cell, then there is some neighborhood $U$ of $e^0_{\alpha}$ such that the only $0$-cell in this neighborhood is $e^0_{\alpha}$; lifting to $ \mathbb{RP}^2 $ gives a neighborhood $U\times D$ for some discrete set $D$ counting the number of sheets; each of these must be homeomorphic to $U$, so that the number of $0$-cells in $ \mathbb{RP}^2 $ that map to $e^0_{\alpha}$ is then $ \left\vert D \right\vert $. Applying the same method to $1$-cells and $2$-cells gives that, for an $n$-sheeted cover $ \mathbb{RP}^2\rightarrow X $, we have that $\chi\left( \mathbb{RP}^2 \right) = n \chi\left( X \right)$. Yet, since $ \chi\left( \mathbb{RP}^2 \right) = 1 $, it follows that $n$ must be $1$ and so must $\chi(X)$. This gives that $X$ is then homeomorphic to $ \mathbb{RP}^2 $ as our map is a $1$-sheeted covering.
\end{solution}
\begin{problem}[Problem 2]
  Let $X$ and $Y$ be subcomplexes of a $2$-dimensional CW complex with finitely many cells. Prove that
  \begin{align*}
    \chi\left( X\cup Y \right) &= \chi\left( X \right) + \chi\left( Y \right) - \chi\left( X\cap Y \right).
  \end{align*}
\end{problem}
\begin{solution}
  We let $V_H, E_H,F_H$ denote the respective number of $0$-cells, $1$-cells, and $2$-cells in the subcomplex $H$. From the principle of inclusion/exclusion, we have that
  \begin{align*}
    V_{X\cup Y} &= V_X + V_Y - V_{X\cap Y}\\
    E_{X\cup Y} &= E_X + E_Y - E_{X\cap Y}\\
    F_{X\cup Y} &= F_X + F_Y - F_{X\cap Y},
  \end{align*}
  which combine to give
  \begin{align*}
    \chi\left( X\cup Y \right) &= \chi\left( X \right) + \chi\left( Y \right) - \chi\left( X\cap Y \right).
  \end{align*}
\end{solution}
\begin{problem}[Problem 3]
  Using their classification, compute the Euler characteristic of all closed surfaces.
\end{problem}
\begin{solution}
  We know that the classification of closed surfaces is given by one of either $S^2$, a $g$-fold connected sum $T^2\#\cdots\# T^2$, or a $g^{\perp}$-fold connected sum $ \mathbb{RP}^2\#\cdots\# \mathbb{RP}^2$.

  Furthermore, we know that the connected sum of any manifolds $M\# N$ has Euler characteristic given by
  \begin{align*}
    \chi\left( M\# N \right) &= \chi\left( M \right) + \chi\left( N \right) - 2.
  \end{align*}
  In particular, this gives that
  \begin{align*}
    \chi\left( S^2 \right) &= 2\\
    \chi\left( T^2\#\cdots\# T^2 \right) &= 2-2g\\
    \chi\left( \mathbb{RP}^2\#\cdots\# \mathbb{RP}^2 \right) &= 1-g^{\perp}.
  \end{align*}
\end{solution}
\begin{problem}[Problem 4]
  Let $M$ be an oriented closed surface of genus $g$ and $p\colon N\rightarrow M$ an $n$-sheeted covering space with $N$ also a closed surface. What is the genus of $N$?
\end{problem}
\begin{solution}
  We see that $M$ has Euler characteristic given by $2-2g$, meaning that $N$ has Euler characteristic given by $2n-2ng$. This gives that $N$ has genus $n\left( g-1 \right) + 1$.
\end{solution}
\end{document}
